\documentclass[11pt]{article}

\usepackage[margin=1in]{geometry}
\usepackage{amsmath}
\usepackage{array}
\usepackage{graphicx}
\usepackage{url}

\title{Computing $x_d$ and $u_d$ from Waypoints with Minimum Snap}
\author{}
\date{}

\begin{document}
\maketitle

\section*{Purpose}

This note summarizes how the current pipeline converts sparse RRT or LLM
waypoints into desired states $x_d$ and feedforward controls $u_d$ using
minimum-snap trajectories.  It also compares the previous separate-clock
scheduling against the newer shared-clock scheduling used for RRT-vs-LLM
calibration scores.

The relevant implementation is in
\[
\text{\path{src/llm_vision_planner/fine_tuning/scripts/min_snap.py}}.
\]

\section*{Step 1: Assign Segment Durations}

Given waypoints
\[
w_0, w_1, \ldots, w_N,
\]
each segment connects $w_i$ to $w_{i+1}$.  The default duration assignment is
\[
T_i = \max\left(\frac{\|w_{i+1}-w_i\|_2}{v_{\mathrm{cruise}}},\ 0.5\right),
\qquad
v_{\mathrm{cruise}} = 0.5\ \mathrm{m/s}.
\]

This is why short LLM segments can get short durations while longer RRT
segments get longer durations.

\section*{Step 2: Fit a Seventh-Degree Polynomial}

For each segment and for each axis, the code uses a seventh-degree polynomial:
\[
p_i(\tau)
=
c_{i,0}
+ c_{i,1}\tau
+ c_{i,2}\tau^2
+ \cdots
+ c_{i,7}\tau^7,
\qquad
0 \leq \tau \leq T_i.
\]

This is done separately for $x(\tau)$ and $y(\tau)$.  The derivatives are
\[
\dot p_i(\tau)
=
c_{i,1}
+2c_{i,2}\tau
+3c_{i,3}\tau^2
+\cdots
+7c_{i,7}\tau^6,
\]
\[
\ddot p_i(\tau)
=
2c_{i,2}
+6c_{i,3}\tau
+12c_{i,4}\tau^2
+\cdots
+42c_{i,7}\tau^5.
\]

In general, the derivative row used by the code is
\[
\frac{d^r}{d\tau^r}\tau^k
=
k(k-1)\cdots(k-r+1)\tau^{k-r}.
\]

\section*{Step 3: Minimum-Snap Optimization}

The trajectory coefficients are selected by minimizing integrated squared snap:
\[
\min_c
\sum_i \int_0^{T_i}
\left(\frac{d^4 p_i(\tau)}{d\tau^4}\right)^2
d\tau.
\]

This becomes a quadratic program:
\[
\min_c \frac{1}{2}c^\top Qc
\quad \text{subject to} \quad Ac=b.
\]

The code solves the equality-constrained KKT system:
\[
\begin{bmatrix}
Q & A^\top\\
A & 0
\end{bmatrix}
\begin{bmatrix}
c\\
\lambda
\end{bmatrix}
=
\begin{bmatrix}
0\\
b
\end{bmatrix}.
\]

The constraints used in the current code are:
\[
p_i(0)=w_i,\qquad p_i(T_i)=w_{i+1},
\]
\[
\dot p_0(0)=0,\qquad \ddot p_0(0)=0,
\]
\[
\dot p_{N-1}(T_{N-1})=0,\qquad
\ddot p_{N-1}(T_{N-1})=0,
\]
and at internal segment boundaries,
\[
\dot p_i(T_i)=\dot p_{i+1}(0),
\qquad
\ddot p_i(T_i)=\ddot p_{i+1}(0).
\]

So the path is smooth in velocity and acceleration at internal waypoints.  The
waypoints are not interpreted as stop points unless they are the first or final
endpoint constraints.

\section*{Step 4: Compute $x_d$ and $u_d$}

At a sample time $t$, the code identifies the active segment and local time
\[
\tau = t - \sum_{j<i} T_j.
\]

Then it evaluates:
\[
p_x(\tau),\quad p_y(\tau),\quad
\dot p_x(\tau),\quad \dot p_y(\tau),\quad
\ddot p_x(\tau),\quad \ddot p_y(\tau).
\]

The desired state is
\[
x_d(t)
=
\begin{bmatrix}
p_x(t)\\
p_y(t)\\
\dot p_x(t)\\
\dot p_y(t)
\end{bmatrix}.
\]

The double-integrator model used by the LQR controller is
\[
\dot p = v,\qquad
\dot v = -d v + d u,
\qquad d=1.1.
\]

To realize a desired acceleration $a_d$, solve
\[
a_d = -d v_d + d u_d.
\]
Therefore,
\[
u_d
=
\frac{a_d+d v_d}{d}
=
v_d + \frac{a_d}{d}.
\]

Per axis, the code computes
\[
u_{x,d} = \frac{\ddot p_x + 1.1\dot p_x}{1.1},
\qquad
u_{y,d} = \frac{\ddot p_y + 1.1\dot p_y}{1.1}.
\]

\section*{Numerical Example: Calibration Sample 1858}

This example uses row \texttt{sample\_id=1858} from
\begin{center}
\path{src/llm_vision_planner/fine_tuning/datasets/conformal_rrt_calibration_dataset_2001.csv}.
\end{center}
The raw LLM output for this row had two waypoints, but the refinement and
verification pipeline produced three verified LLM waypoints.  Minimum snap is
run on the verified waypoints:
\[
\begin{aligned}
\text{RRT}:&\quad
(1.51,2.53),\ (0.83,2.65),\ (0.47,3.07),\\
\text{LLM}:&\quad
(1.51,2.53),\ (0.99,2.63),\ (0.47,3.07).
\end{aligned}
\]

The natural durations are:
\[
\begin{aligned}
T_{\mathrm{RRT}} &= [1.3810,\ 1.1063],\\
T_{\mathrm{LLM,separate}} &= [1.0591,\ 1.3624].
\end{aligned}
\]

Under the shared-clock rule used in the calibration CSV, the LLM uses the RRT
schedule:
\[
T_{\mathrm{LLM,shared}} = [1.3810,\ 1.1063].
\]

\section*{Separate Clock}

With separate clocks, RRT and LLM each receive durations from their own segment
lengths.  At the same wall-clock time $t=0.6$ s:
\[
\text{RRT first-segment progress} = \frac{0.6}{1.3810} = 43.4\%,
\]
\[
\text{LLM first-segment progress} = \frac{0.6}{1.0591} = 56.7\%.
\]

So both are evaluated at the same time, but not at the same local phase of their
first segment.

\begin{center}
\begin{tabular}{>{\raggedright\arraybackslash}p{0.18\linewidth}
                >{\raggedright\arraybackslash}p{0.42\linewidth}
                >{\raggedright\arraybackslash}p{0.28\linewidth}}
\hline
Trajectory & $x_d=[p_x,p_y,v_x,v_y]$ at $t=0.6$ & $u_d=[u_x,u_y]$ at $t=0.6$\\
\hline
RRT & $[1.411,\ 2.519,\ -0.447,\ -0.026]$ & $[-1.484,\ 0.082]$\\
LLM separate & $[1.366,\ 2.540,\ -0.619,\ 0.067]$ & $[-1.820,\ 0.390]$\\
\hline
\end{tabular}
\end{center}

Using the same $dt=0.1$ scoring step as the dataset generation script, the
separate-clock comparison gives
\[
s_u = 0.871456,\qquad s_x = 0.265008.
\]

\section*{Shared Clock}

With the shared-clock rule, RRT is generated normally, and the LLM min-snap
trajectory is generated with the same segment durations as RRT.  The LLM still
uses minimum snap; only the time allocation changes.

At $t=0.6$ s:
\[
\text{RRT first-segment progress}
=
\text{LLM first-segment progress}
=
\frac{0.6}{1.3810}
=
43.4\%.
\]

\begin{center}
\begin{tabular}{>{\raggedright\arraybackslash}p{0.18\linewidth}
                >{\raggedright\arraybackslash}p{0.42\linewidth}
                >{\raggedright\arraybackslash}p{0.28\linewidth}}
\hline
Trajectory & $x_d=[p_x,p_y,v_x,v_y]$ at $t=0.6$ & $u_d=[u_x,u_y]$ at $t=0.6$\\
\hline
RRT & $[1.411,\ 2.519,\ -0.447,\ -0.026]$ & $[-1.484,\ 0.082]$\\
LLM shared & $[1.454,\ 2.514,\ -0.270,\ -0.048]$ & $[-1.011,\ 0.023]$\\
\hline
\end{tabular}
\end{center}

For the calibration CSV row, the shared-clock score is
\[
s_u = 0.836826,\qquad s_x = 0.268212.
\]

\section*{Interpretation}

Separate-clock scoring measures both path deviation and timing deviation:
\[
\text{score} =
\text{path mismatch}
+
\text{schedule mismatch}
+
\text{derivative amplification}.
\]

Shared-clock scoring removes the schedule mismatch:
\[
\text{score} =
\text{path mismatch under a common execution schedule}
+
\text{derivative amplification}.
\]

It does not remove derivative amplification.  Since
\[
u_d = v_d + \frac{a_d}{1.1},
\]
small waypoint differences can still become larger control differences,
especially for short durations.  Shared-clock scheduling only prevents the LLM
from being penalized because its first segment happened to receive a much shorter
duration than the RRT segment.

The number of high-level waypoints does not increase.  However, the number of
sampled trajectory points changes with duration because the implementation uses
\[
\text{samples per segment} = \left\lceil \frac{T_i}{dt} \right\rceil.
\]

\section*{Plots}

\begin{figure}[h]
\centering
\includegraphics[width=0.95\linewidth]{min_snap_waypoint_paths.pdf}
\caption{RRT and LLM waypoints with their generated minimum-snap paths under
separate-clock and shared-clock scheduling.}
\end{figure}

\begin{figure}[h]
\centering
\includegraphics[width=0.95\linewidth]{min_snap_clock_timeseries.pdf}
\caption{Desired $x$ position and feedforward $u_x$ over time.  The separate
LLM clock creates a larger early input spike because its first segment is much
shorter.}
\end{figure}

\end{document}
